\documentclass[notes=show]{beamer}
\usepackage{mathpazo}
\usepackage{hyperref}
\usepackage{multimedia}
\usepackage{tcolorbox}
\usepackage{tikz}
\usetikzlibrary{shadows}
\usetheme{metropolis}
\usecolortheme{default}
\setbeamertemplate{itemize item}{%
    \begin{tikzpicture}
        \shade[ball color=black!100!yellow, preaction={fill=black,
opacity=.00,transform canvas={xshift=1mm,yshift=-1mm, yscale=0.5}}] (0,0) circle (0.6ex);
    \end{tikzpicture}
}
\setbeamertemplate{itemize subitem}{%
    \begin{tikzpicture}
        \shade[ball color=black!100!white, preaction={fill=black,
        opacity=.00,transform canvas={xshift=1mm,yshift=-1mm, yscale=0.5}}] (0,0) circle (0.6ex);
    \end{tikzpicture}
}
\setbeamercolor{frametitle}{bg=white}
\setbeamercolor{frametitle}{fg=black}
\setbeamercolor{background canvas}{bg=white}
\setbeamercolor{block body}{bg=mDarkTeal!30}
\setbeamercolor{block title}{bg=mDarkTeal,fg=black!2}
\setbeamertemplate{navigation symbols}{}
\setbeamertemplate{footline}[page number]
\newenvironment{stepenumerate}{\begin{enumerate}[<+->]}{\end{enumerate}}
\newenvironment{stepitemize}{\begin{itemize}[<+->]}{\end{itemize} }
\newenvironment{stepenumeratewithalert}{\begin{enumerate}[<+-| alert@+>]}{\end{enumerate}}
\newenvironment{stepitemizewithalert}{\begin{itemize}[<+-| alert@+>]}{\end{itemize} }

\begin{document}

\title[Skill, Tasks and Technologies]{Skills, Tasks and Technologies: Implications for Employment and Earnings}
\subtitle{}
\date{Daron Acemoglu \& David Autor \bigskip \\
\textit{Handbook of Labor Economics}, January 2011, Part 2, pages 1118-1166}
\author{}
\maketitle

\begin{frame}{Table of Contents}
\begin{itemize}
\item[\textcolor{gray}{1.}] \textcolor{gray}{Introduction} \medskip
\item[\textcolor{gray}{2.}] \textcolor{gray}{An Overview of Labor Market Trends} \medskip
\item[\textcolor{gray}{3.}] \textcolor{gray}{The Canonical Model} \medskip
\item[\textcolor{black}{4.}] \textcolor{black}{A Ricardian Model of the Labor Market} \medskip
\item[\textcolor{black}{5.}] \textcolor{black}{Comparative Advantange and Wages: An Empirical Approach} \medskip
\item[\textcolor{black}{6.}] \textcolor{black}{Concluding Remarks}
\end{itemize}
\end{frame}

\section{4. A Ricardian Model of the Labor Market}

\subsection{4.1 Environment}

\begin{frame}{4. A Ricardian model of the labor market}
\begin{itemize}
\item[\textcolor{red}{4.1}] \textcolor{red}{Environment}
\item[\textcolor{gray}{4.2}] \textcolor{gray}{Equilibrium without machines}
\item[\textcolor{gray}{4.3}] \textcolor{gray}{Special cases}
\item[\textcolor{gray}{4.4}] \textcolor{gray}{Comparative statics}
\item[\textcolor{gray}{4.5}] \textcolor{gray}{Task replacing technologies}
\item[\textcolor{gray}{4.6}] \textcolor{gray}{Endogenous choice of skill supply}
\item[\textcolor{gray}{4.7}] \textcolor{gray}{Offshoring}
\item[\textcolor{gray}{4.8}] \textcolor{gray}{Directed technical change}
\end{itemize}
\end{frame}

\begin{frame}{Production technology}
\begin{itemize}
\item Static environment with a unique final good, $Y$. \medskip
\item $Y$ produced with a continuum of tasks on the interval $[0,1]$. \medskip
\item Cobb-Douglas technology mapping tasks into the final good:
\[
Y= \exp \left[ \int_{0}^{1} \ln y(i)di \right] \tag{11}  \label{eq11}
\]
where $y(i)$ is the service or production of task $i$. \medskip
\item The final good is the numeraire, $P \equiv 1$.
\end{itemize}
\end{frame}

\begin{frame}{Supply of skills to tasks}
\begin{itemize}
\item There are 3 types of labor: low, medium and high skilled. Their supply is inelastic and given by $L, M, H$. \medskip
\item Each task on the continuum is produced using:
\begin{align*}
y(i) & = A_{L} \alpha_{L}(i) l(i) + A_{M} \alpha_{M}(i) m(i) \\
 & + A_{H} \alpha_{H}(i) h(i) + A_{K} \alpha_{K}(i) k(i) \tag{12}  \label{eq12}
\end{align*}
where:
\begin{itemize}
\item $A_{L}, A_{M}, A_{H}$ are skill productivity parameters. \medskip
\item $ \alpha_{L}(i), \alpha_{M}(i), \alpha_{H}(i)$ are task productivity schedules determining comparative advantage in task production. \medskip
\item $ l(i), m(i), h(i)$ are the number of workers allocated to task $i$.
\end{itemize}
\end{itemize}
\end{frame}

\begin{frame}{Comparative advantage in task production}
\begin{block}{Assumption 1}
$ \alpha_{L}(i)/\alpha_{M}(i)$ and $ \alpha_{M}(i)/\alpha_{H}(i)$ are continuously differentiable and strictly decreasing in $i$.
\end{block}
\begin{itemize}
\item This is an assumption about the \textit{relative} productivity of workers in doing different tasks. \medskip
\item This assumption results in the allocation of workers across tasks based on their comparative advantage (as in Roy [51]).\medskip
\item The rationale is that higher $i$ correspond to ``more complex'' tasks for which more skilled workers have a comparative advantage.
\end{itemize}
\end{frame}

\begin{frame}{Clearing factor markets}
\begin{itemize}
\item Labor markets clearing requires:
\[
\int_{0}^{1} l(i)di = L \text{ and } \int_{0}^{1} m(i)di = M \text{ and } \int_{0}^{1} h(i)di = H \tag{13}  \label{eq13}
\]
\item The role of capital as a factor of production: \medskip
\begin{itemize}
\item \underline{Equilibrium without machines} (sections 4.2-4.4): \\
$ \alpha_{K}(i) = 0$ for all $i$. \medskip
\item \underline{Task replacing technologies} (section 4.5): \\
$ \alpha_{K}(i) \neq 0$ for some $i$.
\end{itemize}
\end{itemize}
\end{frame}

\subsection{4.2 Equilibrium without machines}

\begin{frame}{4. A Ricardian model of the labor market}
\begin{itemize}
\item[\textcolor{gray}{4.1}] \textcolor{gray}{Environment}
\item[\textcolor{red}{4.2}] \textcolor{red}{Equilibrium without machines}
\item[\textcolor{gray}{4.3}] \textcolor{gray}{Special cases}
\item[\textcolor{gray}{4.4}] \textcolor{gray}{Comparative statics}
\item[\textcolor{gray}{4.5}] \textcolor{gray}{Task replacing technologies}
\item[\textcolor{gray}{4.6}] \textcolor{gray}{Endogenous choice of skill supply}
\item[\textcolor{gray}{4.7}] \textcolor{gray}{Offshoring}
\item[\textcolor{gray}{4.8}] \textcolor{gray}{Directed technical change}
\end{itemize}
\end{frame}

\begin{frame}{Task thresholds}
\begin{block}{Lemma 1}
In any equilibrium there exist $I_{L}$ and $I_{H}$ such that $0 < I_{L} < I_{H} < 1$ and for any $i < I_{L}$ we have that $m(i)=h(i)=0$, for any $i \in (I_{L},I_{H})$ we have that $l(i)=h(i)=0$, and for any $i > I_{H}$ we have that $l(i)=m(i)=0$.
\end{block}
Thresholds $I_{L}, I_{H}$ determine the allocation of tasks to skill: \medskip
\begin{itemize}
\item $i < I_{L}$ will be performed by $L$ workers (Manual). \medskip
\item $I_{L} \leq i \leq I_{H}$ will be performed by $M$ workers (Routine). \medskip
\item $i > I_{H}$ will be performed by $H$ workers (Abstract).
\end{itemize}
\end{frame}

\begin{frame}{Task thresholds $I_{L}$ and $I_{H}$ must exist}
\begin{itemize}
\item Given that $\alpha_{L}(i) / \alpha_{M}(i)$ is strictly decreasing in $i$, there must exist a task $I_{L}$ such that the unit cost of producing $I_{L}$ is the same for $L$ and $M$ workers:
\[
\frac{w_{L}}{A_{L} \alpha_{L}(I_{L})} = \frac{w_{M}}{A_{M} \alpha_{M}(I_{L})}
\]
\item Given that $\alpha_{M}(i) / \alpha_{H}(i)$ is strictly decreasing in $i$, there must exist a task $I_{H}$ such that the unit cost of producing $I_{H}$ is the same for $M$ and $H$ workers:
\[
\frac{w_{M}}{A_{M} \alpha_{M}(I_{H})} = \frac{w_{H}}{A_{M} \alpha_{M}(I_{H})}
\]
\end{itemize}
\end{frame}

\begin{frame}{Three equilibrium conditions}
\begin{enumerate}
\item[E1.] Law of one price for skill. \bigskip
\item[E2.] Equal task cost shares.\bigskip
\item[E3.] No arbitrage across skills at task thresholds. \bigskip
\item[E4.] Price of the final good equals marginal cost.
\end{enumerate}
\end{frame}

\begin{frame}{E1. Law of one price for skill}
\begin{itemize}
\item Denote $p(i)$ as the price of task $i$. \medskip
\item All tasks by $j=L,M,H$ workers must pay the same wage $w_{j}$:
\begin{tcolorbox}
\[
w_{j} = p(i)A_{j} \alpha_{j}(i) \text{ for } j = L,M,H
\]
\end{tcolorbox}
\item This has a convenient implication:
\[
p(i) \alpha_{L}(i) = p(i') \alpha_{L}(i') \equiv P_{L} \text{ for any } i,i' < I_{L} \tag{14} \label{eq14}
\]
\[
p(i) \alpha_{M}(i) = p(i') \alpha_{M}(i') \equiv P_{M} \text{ for any } I_{L} < i,i' < I_{H} \tag{15} \label{eq15}
\]
\[
p(i) \alpha_{H}(i) = p(i') \alpha_{H}(i') \equiv P_{H} \text{ for any } i > I_{H} \tag{16} \label{eq16}
\]
\end{itemize}
\end{frame}

\begin{frame}{E2. Equal task cost shares within skill groups}
\begin{itemize}
\item Cobb-Douglas technology implies:
\[
p(i)y(i) = p(i')y(i') \text{ for any } i,i' \in [0,1] \tag{17} \label{eq17}
\]
\item Within the group of $L$ workers this implies:
\[
p(i) \alpha_{L}(i) l(i) = p(i') \alpha_{L}(i') l(i') \text{ for any } i,i' < I_{L}
\]
\item From the law of one price for skill we had that:
\[
p(i) \alpha_{L}(i) = p(i') \alpha_{L}(i') \equiv P_{L}
\]
\item This implies:
\[
l(i) = l(i') \text{ for any } i,i' < I_{L}
\]
\end{itemize}
\end{frame}

\begin{frame}{E2. Equal task cost shares within skill group}
\begin{itemize}
\item Within skill groups, we had that:
\[
l(i) = l(i') \text{ for any } i,i' < I_{L}
\]
\[
m(i) = m(i') \text{ for any } I_{L} < i,i' < I_{H}
\]
\[
h(i) = h(i') \text{ for any } i,i' > I_{H}
\]
\item Using the market clearing conditions, this gives:
\begin{tcolorbox}
\[
l(i) = \frac{L}{I_{L}} \text{ for any } i < I_{L} \tag{18} \label{eq18}
\]
\[
m(i) = \frac{M}{I_{H} - I_{L}} \text{ for any } I_{L} < i < I_{H} \tag{19} \label{eq19}
\]
\[
h(i) = \frac{H}{1 - I_{H}} \text{ for any } i > I_{H} \tag{20} \label{eq20}
\]
\end{tcolorbox}
\end{itemize}
\end{frame}

\begin{frame}{E2. Equal task cost shares between skill group}
\begin{itemize}
\item Cobb-Douglas technology implies:
\[
p(i)y(i) = p(i')y(i') \text{ for any } i,i' \in [0,1] \tag{17} \label{eq17}
\]
\item Between $M$ and $H$ workers this implies:
\[
p(i) \alpha_{M}(i) A_{M} m(i) = p(i') \alpha_{H}(i') A_{H} h(i')
\]
\item Using definitions and results from above:
\[
\frac{P_{M}A_{M}M}{I_{H} - I_{L}} = \frac{P_{H}A_{H}H}{1 - I_{H}}
\]
\[
\frac{P_{H}}{P_{M}} = \left[\frac{A_{H}H}{1 - I_{H}}\right] ^{-1} \left[\frac{A_{M}M}{I_{H} - I_{L}} \right] \tag{21} \label{eq21}
\]
\end{itemize}
\end{frame}

\begin{frame}{E2. Equal task cost shares between skill group}
\begin{itemize}
\item We had that:
\[
\frac{P_{H}}{P_{M}} = \left[\frac{A_{H}H}{1 - I_{H}}\right] ^{-1} \left[\frac{A_{M}M}{I_{H} - I_{L}} \right] \tag{21} \label{eq21}
\]
\item Similarly, comparing tasks between $L$ and $M$ workers gives:
\[
\frac{P_{M}}{P_{L}} = \left[\frac{A_{M}M}{I_{H} - I_{L}}\right] ^{-1} \left[\frac{A_{L}L}{I_{L}} \right] \tag{22} \label{eq22}
\]
\item We can use these equations to solve for relative wages as a function of task thresholds.
\end{itemize}
\end{frame}

\begin{frame}{E2. Equal task cost shares between skill group}
\begin{itemize}
\item Using eq. (21), we get that:
\begin{tcolorbox}
\[
\frac{w_{H}}{w_{M}} = \frac{P_{H}A_{H}}{P_{M}A_{M}} = \left[\frac{1-I_{H}}{I_{H}-I_{L}} \right] \left[ \frac{M}{H} \right] \tag{26} \label{eq26}
\]
\end{tcolorbox}
\item Using eq. (22), we get that:
\begin{tcolorbox}
\[
\frac{w_{M}}{w_{L}} = \frac{P_{M}A_{M}}{P_{L}A_{L}} = \left[\frac{I_{H}-I_{L}}{I_{L}} \right] \left[ \frac{L}{M} \right] \tag{27} \label{eq27}
\]
\end{tcolorbox}
\item Given $I_{L}$ and $I_{H}$, these equations solve for relative wages.
\end{itemize}
\end{frame}

\begin{frame}{E3. No arbitrage across skills at task thresholds}
\begin{itemize}
\item Unit cost to produce $I_{H}$ must be the same for $M$ or $H$:
\[
\frac{w_{H}}{w_{M}} = \frac{A_{H} \alpha_{H}(I_{H})}{A_{M} \alpha_{M}(I_{H})} = \left[\frac{1-I_{H}}{I_{H}-I_{L}} \right] \left[ \frac{M}{H} \right] \tag{23} \label{eq23}
\]
\item Unit cost to produce $I_{L}$ must be the same for $L$ or $M$:
\[
\frac{w_{M}}{w_{L}}  = \frac{A_{M} \alpha_{M}(I_{L})}{A_{L} \alpha_{L}(I_{L})} = \left[\frac{I_{H}-I_{L}}{I_{L}} \right] \left[ \frac{L}{M} \right]  \tag{24} \label{eq24}
\]
\item These equations solve for $I_{L}$ and $I_{H}$. \medskip
\item This solves for $l(i), m(i), h(i)$, $P_{H}/P_{M}, P_{M}/P_{L}$ and $w_{H}/w_{M}, w_{M}/w_{L}$.
\end{itemize}
\end{frame}

\begin{frame}{Summary of equilibrium}
\begin{block}{Proposition 1}
There exists a unique equilibrium summarized by $(I_{L}, I_{H}, P_{L}, P_{M}, P_{H}, w_{L}, w_{M}, w_{H})$.
\end{block}
\begin{itemize}
\item There exist values of $I_{L}$ and $I_{H}$. Still required is proof that these values are unique. \medskip
\item Knowing $I_{L}$ and $I_{H}$ directly solves for $l(i), m(i), h(i)$, $P_{H}/P_{M}, P_{M}/P_{L}$ and $w_{H}/w_{M}, w_{M}/w_{L}$. \medskip
\item Solving for $P_{H}, P_{M}, P_{L}$ and $w_{H}, w_{M}, w_{L}$ requires solving for the equilibrium price of the final good $P$.
\end{itemize}
\end{frame}

\begin{frame}{Uniqueness of $I_{L}$ and $I_{H}$}
\begin{itemize}
\item Rewrite the no arbitrage condition for $I_{H}$ as:
\[
\frac{A_{M} \alpha_{M}(I_{H})M}{A_{H} \alpha_{H}(I_{H})H} \left[ \frac{1 - I_{H}}{I_{H}-I_{L}} \right] =1 \tag{23'} \label{eq23'}
\]
\item Rewrite the no arbitrage condition for $I_{L}$ as:
\[
\frac{A_{L} \alpha_{L}(I_{L})L}{A_{M} \alpha_{M}(I_{L})M} \left[ \frac{I_{H} - I_{L}}{I_{L}} \right] = 1 \tag{24'} \label{eq24'}
\]
\item Both are upward sloping in ($I_{L}, I_{H}$)-space but eq. (23') is everywhere steeper than eq. (24') given Assumption 1.
\end{itemize}
\end{frame}

\newpage
\begin{center}
\begin{figure}
\includegraphics[height=8.5cm,width=\textwidth]{figure22.png}
\end{figure} 
\end{center}
\newpage

\begin{frame}{Equilibrium allocation of skills to tasks}
\begin{itemize}
\item Rewrite the no arbitrage condition for $I_{H}$ as:
\[
\frac{A_{H}H}{A_{M}M} = \frac{1-I_{H}}{I_{H} - I_{L}} \frac{\alpha_{M}(I_{H})}{\alpha_{H}(I_{H})}  \tag{29} \label{eq29}
\]
which is decreasing in $I_{H}$ for given $I_{L}$. \medskip
\item Rewrite the no arbitrage condition for $I_{L}$ as:
\[
\frac{A_{M}M}{A_{L}L} = \frac{I_{H} - I_{L}}{I_{L}} \frac{\alpha_{L}(I_{L})}{\alpha_{M}(I_{L})} 
\]
which is decreasing in $I_{L}$ for given $I_{H}$.
\end{itemize}
\end{frame}

\newpage
\begin{center}
\begin{figure}
\includegraphics[height=8.5cm,width=\textwidth]{figure23.png}
\end{figure} 
\end{center}
\newpage

\begin{frame}{E.4 Price of the final good equals marginal cost}
\begin{itemize}
\item Profit maximization implies:
\[
P = MC= \exp \left[ \int_{0}^{1} \ln(p(i))di \right]  \equiv 1
\]
\item Rewriting gives:
\begin{align*}
\ln P & = \int_{0}^{I_{L}} [\ln(P_{L}) - \ln(\alpha_{L}(i))]di + \int_{I_{L}}^{I_{H}} [\ln(P_{M}) - \ln(\alpha_{M}(i))]di \\
& + \int_{I_{H}}^{1} [\ln(P_{H}) - \ln(\alpha_{H}(i))]di= 0 \tag{28} \label{eqs28}
\end{align*}
which, together with eqs. (21)-(22), solves for $ P_{L}, P_{M}, P_{H} $. \medskip
\item This solves for $w_{j}=P_{j}A_{j} $ for $ j=L,M,H $.
\end{itemize}
\end{frame}

\subsection{4.3 Special cases}

\begin{frame}{4. A Ricardian model of the labor market}
\begin{itemize}
\item[\textcolor{gray}{4.1}] \textcolor{gray}{Environment}
\item[\textcolor{gray}{4.2}] \textcolor{gray}{Equilibrium without machines}
\item[\textcolor{red}{4.3}] \textcolor{red}{Special cases}
\item[\textcolor{gray}{4.4}] \textcolor{gray}{Comparative statics}
\item[\textcolor{gray}{4.5}] \textcolor{gray}{Task replacing technologies}
\item[\textcolor{gray}{4.6}] \textcolor{gray}{Endogenous choice of skill supply}
\item[\textcolor{gray}{4.7}] \textcolor{gray}{Offshoring}
\item[\textcolor{gray}{4.8}] \textcolor{gray}{Directed technical change}
\end{itemize}
\end{frame}

\begin{frame}{Special cases}
\begin{itemize}
\item We discuss 2 special cases: \medskip
\begin{enumerate} \medskip
\item A model with a fixed task set for L-type workers. \medskip
\item Assuming two skill groups and simple versions of comparative advantage schedules (as in Acemoglu \& Zilibotti [01]).\medskip
\end{enumerate}
\item There are other interesting cases: \medskip
\begin{itemize}
\item Assuming a discrete set of tasks instead and a continuum of skills (as in Autor, Levy \& Murnane [03]). \medskip
\item Assuming CES instead of Cobb-Douglas task production function (as in Acemoglu \& Restrepo [18,19])
\end{itemize}
\end{itemize}
\end{frame}

\begin{frame}{Special case 1: Fixed task set for L-type workers}
\begin{itemize}
\item Assume 3 types of workers: $L,M,H$. \medskip
\item Assume the following comparative advantage schedule for low skill workers:
\[
\alpha_{L}(i) = 
\begin{cases}
\tilde{\alpha}_{L} \text{ if } i \leq \tilde{I}_{L} \\
0 \text{ if } i > \tilde{I}_{L} \\
\end{cases}
\]
where $\tilde{\alpha}_{L} $ is large and $ \tilde{I}_{L}$ is small. \medskip
\item This fixes  $I_{L} = \tilde{I}_{L}$ while $I_{H}$ is determined as before from an arbitrage condition. \medskip
\item This is essentially a two-skill version of the model.  
\end{itemize}
\end{frame}

\newpage
\begin{center}
\begin{figure}
\includegraphics[height=8.5cm,width=\textwidth]{figure24.png}
\end{figure} 
\end{center}
\newpage

\begin{frame}{Special case 2: Assuming two skill groups}
\begin{itemize}
\item Assume there are only 2 types of workers, low and high skill. \medskip
\item Assume the following comparative advantage schedules:
\[
\alpha_{L}(i) = 1-i \text{ and } \alpha_{H}(i)=i
\]
\item This model is isomorphic to the canonical model with $\sigma =2$:
\begin{align*}
\omega \equiv \frac{w_{H}}{w_{L}} & = \left[ \frac{A_{H}}{A_{L}} \right] ^{ \frac{ 1}{2}}  \left[ \frac{H}{L} \right] ^ {-\frac{1}{ 2}} \tag{5} \label{eq5} \\ 
& = \left[ \frac{A_{H}}{A_{L}} \right] ^{ \frac{ \sigma -1}{ \sigma}}  \left[ \frac{H}{L} \right] ^ {-\frac{1}{ \sigma}}  
\end{align*}
\end{itemize}
\end{frame}

\begin{frame}{Special case 2: Assuming two skill groups}
\begin{itemize}
\item Equal division of labor among tasks within skill groups gives:
\[
l(i) = \frac{L}{I} \text{ and } h(i) = \frac{H}{1-I}
\]
\item Comparing tasks between skill groups gives:
\[
\frac{P_{H}}{P_{L}} = \left[ \frac{A_{H}H}{1-I} \right]^{-1} \left[ \frac{A_{L}L}{I} \right]
\]
\item The relative wage is given by:
\[
\frac{w_{H}}{w_{L}} = \frac{P_{H}A_{H}}{P_{L}A_{L}} = \left[ \frac{H}{1-I} \right]^{-1} \left[ \frac{L}{I} \right]
\]
\end{itemize}
\end{frame}

\begin{frame}{Special case 2: Assuming two skill groups}
\begin{itemize}
\item No arbitrage across skill gives:
\[
\frac{A_{L} \alpha_{L}(I)L}{I} = \frac{A_{H} \alpha_{H}(I)H}{1-I} 
\]
\item Re-writing and using that $\alpha_{L}(I)=1-I$ and $\alpha_{H}(I)=I$ gives:
\[
\frac{1-I}{I} = \left[ \frac{A_{H}H}{A_{L}L} \right]^{\frac{1}{2}}
\]
\item Re-writing the relative wage gives using this expression gives:
\[
\frac{w_{H}}{w_{L}} =  \left[ \frac{A_{H}}{A_{L}} \right]^{\frac{1}{2}}  \left[ \frac{H}{L} \right]^{-\frac{1}{2}}
\]
\end{itemize}
\end{frame}

\subsection{4.4 Comparative statics}

\begin{frame}{4. A Ricardian model of the labor market}
\begin{itemize}
\item[\textcolor{gray}{4.1}] \textcolor{gray}{Environment}
\item[\textcolor{gray}{4.2}] \textcolor{gray}{Equilibrium without machines}
\item[\textcolor{gray}{4.3}] \textcolor{gray}{Special cases}
\item[\textcolor{red}{4.4}] \textcolor{red}{Comparative statics}
\item[\textcolor{gray}{4.5}] \textcolor{gray}{Task replacing technologies}
\item[\textcolor{gray}{4.6}] \textcolor{gray}{Endogenous choice of skill supply}
\item[\textcolor{gray}{4.7}] \textcolor{gray}{Offshoring}
\item[\textcolor{gray}{4.8}] \textcolor{gray}{Directed technical change}
\end{itemize}
\end{frame}

\begin{frame}{Basic comparative statics}
\begin{itemize}
\item Taking logs of eq. (23') gives:
\begin{align*}
\ln (A_{M}/A_{H}) + \beta_{H}(I_{H}) + \ln (M/H) \\
+ \ln ([1 - I_{H}]/[I_{H} - I_{L}]) = 0  \tag{32} \label{eq32}
\end{align*}
with $ \beta_{H}(I_{H}) \equiv \ln ( \alpha_{M}(I_{H}) / \alpha_{H}(I_{H}))$. \medskip 
\item Taking logs of eq. (24') gives:
\begin{align*}
\ln (A_{L}/A_{M}) + \beta_{L}(I_{L}) + \ln (L/M) \\
+ \ln ([I_{H} - I_{L}]/I_{L}) = 0 \tag{33} \label{eq33}
\end{align*}
with $ \beta_{L}(I_{L}) \equiv \ln ( \alpha_{L}(I_{L}) / \alpha_{M}(I_{L}))$.
\end{itemize}
\end{frame}

\newpage
\begin{center}
\begin{figure}
\includegraphics[height=8.5cm,width=\textwidth]{figure25.png}
\end{figure} 
\end{center}
\newpage

\newpage
\begin{center}
\begin{figure}
\includegraphics[height=8.5cm,width=\textwidth]{figure26.png}
\end{figure} 
\end{center}
\newpage

\begin{frame}{Proposition 2.1: The response of task allocation to $A_{H},H$}
\begin{enumerate}
\item Task range of $H$ workers:
\[
\frac{dI_{H}}{d \ln (A_{H})} = \frac{d I_{H}}{d \ln (H)} < 0
\]
\item Task range of $L$ workers:
\[
\frac{dI_{L}}{d \ln (A_{H})} = \frac{d I_{L}}{d \ln (H)} < 0
\]
\item Task range for $M$ workers:
\[
\frac{d (I_{H} - I_{L})}{d \ln (A_{H})} = \frac{d (I_{H} - I_{L})}{d \ln (H)} < 0
\]
\end{enumerate}
\end{frame}

\begin{frame}{Proposition 2.1: The response of task allocation to $A_{L},L$}
\begin{enumerate}
\item Task range of $H$ workers:
\[
\frac{dI_{H}}{d \ln (A_{L})} = \frac{d I_{H}}{d \ln (L)} > 0
\]
\item Task range of $L$ workers:
\[
\frac{dI_{L}}{d \ln (A_{L})} = \frac{d I_{L}}{d \ln (L)} > 0
\]
\item Task range for $M$ workers:
\[
\frac{d (I_{H} - I_{L})}{d \ln (A_{L})} = \frac{d (I_{H} - I_{L})}{d \ln (L)} < 0
\]
\end{enumerate}
\end{frame}

\begin{frame}{Proposition 2.1: The response of task allocation to $A_{M},M$}
\begin{enumerate}
\item Task range of $H$ workers:
\[
\frac{dI_{H}}{d \ln (A_{M})} = \frac{d I_{H}}{d \ln (M)} > 0
\]
\item Task range of $L$ workers:
\[
\frac{dI_{L}}{d \ln (A_{M})} = \frac{d I_{L}}{d \ln (M)} < 0
\]
\item Task range for $M$ workers:
\[
\frac{d (I_{H} - I_{L})}{d \ln (A_{M})} = \frac{d (I_{H} - I_{L})}{d \ln (M)} > 0
\]
\end{enumerate}
\end{frame}

\begin{frame}{Proposition 2.2: Response of relative wages to skill supplies}
\begin{enumerate}
\item An increase in $H$:
\[
\frac{d \ln (w_{H}/w_{M})}{d \ln (H)} < 0 \text{ and } \frac{d \ln (w_{H}/w_{L})}{d \ln (H)} < 0
\]
\item An increase in $M$:
\[
\frac{d \ln (w_{H}/w_{M})}{d \ln (M)} > 0 \text{ and } \frac{d \ln (w_{M}/w_{L})}{d \ln (M)} < 0
\]
\item An increase in $L$:
\[
\frac{d \ln (w_{M}/w_{L})}{d \ln (L)} > 0  \text{ and } \frac{d \ln (w_{H}/w_{L})}{d \ln (L)} > 0
\]
\end{enumerate}
\end{frame}

\begin{frame}{Proposition 2.3: Response of relative wages to $A_{H}, A_{M}, A_{L}$}
\begin{enumerate}
\item An increase in $A_{H}$:
\[
\frac{d \ln (w_{H}/w_{M})}{d \ln (A_{H})} > 0  \text{ and } \frac{d \ln (w_{H}/w_{L})}{d \ln (A_{H})} > 0
\]
\item An increase in $A_{M}$:
\[
\frac{d \ln (w_{H}/w_{M})}{d \ln (A_{M})} < 0 \text{ and } \frac{d \ln (w_{M}/w_{L})}{d \ln (A_{M})} > 0
\]
\item An increase in $A_{L}$:
\[
\frac{d \ln (w_{M}/w_{L})}{d \ln (A_{L})} < 0 \text{ and } \frac{d \ln (w_{H}/w_{L})}{d \ln (A_{L})} < 0
\]
\end{enumerate}
\end{frame}

\begin{frame}{Propositions 2.2 and 2.3: Response of $w_{H}/w_{L}$ to $M,A_{M}$}
\begin{itemize}
\item This depends critically on the comparative advantages of marginal $H$ and $L$ workers in doing $M$ tasks, given by:
\[
\beta_{H}(I_{H}) \equiv \ln ( \alpha_{M}(I_{H}) / \alpha_{H}(I_{H}))
\]
\[
\beta_{L}(I_{L}) \equiv \ln ( \alpha_{L}(I_{L}) / \alpha_{M}(I_{L}))
\]
\item Define:
\[
\beta_{L}'(I_{L}) I_{L} \equiv \partial \beta_{L}(I_{L})/ \partial I_{L}  
\]
\[
\beta_{H}'(I_{H}) I_{H} \equiv \partial \beta_{H}(I_{H})/ \partial I_{H}  
\]
\item If $ \beta_{H}'(I_{H})$ is high relative to $\beta_{L}'(I_{L})$ in absolute value, $H$ workers have a relatively strong comparative advantage for tasks above $I_{H}$ compared to $L$ workers for tasks below $I_{L}$.
\end{itemize}
\end{frame}

\begin{frame}{Propositions 2.2 and 2.3: Response of $w_{H}/w_{L}$ to $M,A_{M}$}
\begin{itemize}
\item An increase in $M$: 
\[
\frac{d \ln (w_{H}/w_{L})}{d \ln (M)} \gtrless 0 \text{ iff } | \beta_{H}'(I_{H}) (1 - I_{H})| \gtrless |\beta_{L}'(I_{L}) I_{L}|
\]
\item An increase in $A_{M}$:
\[
\frac{d \ln (w_{H}/w_{L})}{d \ln (A_{M})} \gtrless 0 \text{ iff } | \beta_{H}'(I_{H}) (1 - I_{H})| \gtrless |\beta_{L}'(I_{L}) I_{L}|
\]
\item If $ \beta_{H}'(I_{H})$ is high relative to $\beta_{L}'(I_{L})$ in absolute value, an increase in $M$ or $A_{M}$ displaces $L$ workers more than $H$ workers (and $I_{L}$ falls more than $I_{H}$ rises) such that $w_{H}/w_{L}$ rises.
\end{itemize}
\end{frame}

\subsection{4.5 Task replacing technologies}

\begin{frame}{4. A Ricardian model of the labor market}
\begin{itemize}
\item[\textcolor{gray}{4.1}] \textcolor{gray}{Environment}
\item[\textcolor{gray}{4.2}] \textcolor{gray}{Equilibrium without machines}
\item[\textcolor{gray}{4.3}] \textcolor{gray}{Special cases}
\item[\textcolor{gray}{4.4}] \textcolor{gray}{Comparative statics}
\item[\textcolor{red}{4.5}] \textcolor{red}{Task replacing technologies}
\item[\textcolor{gray}{4.6}] \textcolor{gray}{Endogenous choice of skill supply}
\item[\textcolor{gray}{4.7}] \textcolor{gray}{Offshoring}
\item[\textcolor{gray}{4.8}] \textcolor{gray}{Directed technical change}
\end{itemize}
\end{frame}

\begin{frame}{Automation of routine labor tasks done by $M$-type workers}
\begin{itemize}
\item Labor tasks most subject to machine displacement are routine or codifiable. \medskip
\item These tasks are primarily done by $M$-type workers. \medskip
\item Assume a range of tasks $[I',I''] \subset [I_{L},I_{H}]$ for which $\alpha_{K}(i)$ increases sufficiently (with fixed cost of capital $r$) so that they are now more economically performed by machines. \medskip
\item For all $i \notin [I',I'']$, continue to assume that $\alpha_{K}(i)=0$. 
\end{itemize}
\end{frame}

\begin{frame}{Changes in task thresholds}
\begin{block}{Proposition 3}
Suppose we start with an equilibrium characterized by thresholds $[I_{L},I_{H}]$ and technical change implies that the tasks in the range $[I',I''] \subset [I_{L},I_{H}]$ are now performed by machines. \\ \medskip
Then: \\ \medskip
After the introduction of machines, there exists a new equilibrium characterized by new thresholds $\hat{I}_{L}$ and $\hat{I}_{H}$ such that $0 < \hat{I}_{L} < I' < I'' < \hat{I}_{H} < 1$ and for any $i < \hat{I}_{L}$ we have that $m(i)=h(i)=0$ and $l(i)=L/\hat{I}_{L}$; for any $i \in (\hat{I}_{L},I') \cup (I'',\hat{I}_{H})$, $l(i)=h(i)=0$ and $m(i)=M/(\hat{I}_{H} - I'' + I' - \hat{I}_{L})$; and for any $i > \hat{I}_{H}$,  $l(i)=m(i)=0$ and $h(i)=H/(1-\hat{I}_{H})$.
\end{block}
\end{frame}

\begin{frame}{Changes in relative wages}
\begin{block}{Proposition 4}
Suppose we start with an equilibrium characterized by thresholds $[I_{L},I_{H}]$ and technical change implies that the tasks in the range $[I',I''] \subset [I_{L},I_{H}]$ are now performed by machines. \\ \medskip
Then:
\begin{enumerate}
\item $w_{H}/w_{M}$ increases; \medskip
\item $w_{M}/w_{L}$ decreases; \medskip
\item $w_{H}/w_{L} \gtrless 0 $ if and only if $ | \beta_{H}'(I_{H}) (1 - I_{H})| \gtrless |\beta_{L}'(I_{L}) I_{L}| $.
\end{enumerate}
\end{block}
\end{frame}

\begin{frame}{Automation of routine labor tasks done by $M$-type workers}
\begin{itemize}
\item Focal case: \smallskip
\begin{itemize}
\item Routine tasks are of average complexity for humans. \smallskip
\item Strong comparative advantage of $H$ relative to $L$ workers. \smallskip
\end{itemize}
\item If so, there is wage and employment polarization: \smallskip
\begin{enumerate}
\item Wages: \smallskip
\begin{itemize}
\item Middle wages fall relative to top and bottom. \smallskip
\item Top rises relative to bottom.
\end{itemize}
\item Employment: \smallskip
\begin{itemize}
\item Declining labor input in middling routine tasks. \smallskip
\item Middle-skill workers move disproportionately downward.
\end{itemize}
\end{enumerate}
\end{itemize}
\end{frame}

\begin{frame}{Ricardian model: Summary}
\begin{itemize}
\item Model's inputs: \smallskip
\begin{enumerate}
\item Explicit distinction between skills and tasks. \smallskip
\item Comparative advantage among workers in different tasks. \smallskip
\item Multiple sources of competing task supplies. \smallskip
\end{enumerate}
\item What the model delivers: \smallskip
\begin{enumerate}
\item A natural concept of occupations. \smallskip
\item Sorting of skill to tasks based on comparative advantage. \smallskip
\item Reallocation of skill across tasks as technology changes.\smallskip
\item Polarization of wages and employment.\smallskip
\item Technological progress can decrease wage levels.
\end{enumerate}
\end{itemize}
\end{frame}

\subsection{4.6 Endogenous choice of skill supply}

\begin{frame}{4. A Ricardian model of the labor market}
\begin{itemize}
\item[\textcolor{gray}{4.1}] \textcolor{gray}{Environment}
\item[\textcolor{gray}{4.2}] \textcolor{gray}{Equilibrium without machines}
\item[\textcolor{gray}{4.3}] \textcolor{gray}{Special cases}
\item[\textcolor{gray}{4.4}] \textcolor{gray}{Comparative statics}
\item[\textcolor{gray}{4.5}] \textcolor{gray}{Task replacing technologies}
\item[\textcolor{red}{4.6}] \textcolor{red}{Endogenous choice of skill supply}
\item[\textcolor{gray}{4.7}] \textcolor{gray}{Offshoring}
\item[\textcolor{gray}{4.8}] \textcolor{gray}{Directed technical change}
\end{itemize}
\end{frame}

\begin{frame}{Endogenous choice of skill supply}
\begin{itemize}
\item We have focussed on the substitution of skills across tasks. \medskip
\item A complementary force is substitution of workers across skills. \medskip
\item In response to changes in technology or factor supplies, workers may change the types of skills they supply. \medskip
\item If machines replace medium skill workers, the supply to especially low skills will increase. \medskip
\item This gives a richer explanation of observed wage and job polarization.
\end{itemize}
\end{frame}

\begin{frame}{Environment}
\begin{itemize}
\item Each worker $j$ is endowed with some amount of low skill $l^{j}$, medium skill $m^{j}$ and high skill $h^{j}$. \medskip
\item Workers have one unit of time subject to a skill allocation constraint:
\[
t_{l}^{j} + t_{m}^{j} + t_{h}^{j} \leqslant 1 
\]
\item The worker's income is:
\[
w_{L}t_{l}^{j}l_{j} + w_{M}t_{m}^{j}m_{j} + w_{H}t_{h}^{j}h_{j} 
\]
\item A worker will prefer to allocate all her time to one skill.
\end{itemize}
\end{frame}

\begin{frame}{Supply to skills}
\begin{itemize}
\item The supply to skills is given by:
\[
L = \int_{j \in E_{l}} l^{j}dj \text{  ;  } M = \int_{j \in E_{m}} m^{j}dj \text{  ;  } H = \int_{j \in E_{h}} h^{j}dj
\]
with $E_{l}, E_{m}, E_{h}$ the sets of workers choosing to supply their labor to low, medium and high skills respectively.
\item A worker will choose to be in set $E_{h}$ only if:
\[
\frac{l^{j}}{h^{j}} \leq \frac{w_{H}}{w_{L}} \text{ and } \frac{m^{j}}{h^{j}} \leq \frac{w_{H}}{w_{M}}
\]
with similar inequalities for sets $E_{m}$ and $E_{l}$.
\end{itemize}
\end{frame}

\begin{frame}{Comparative advantage in skill}
\begin{block}{Assumption 2}
$h^{j}/m^{j}$ and $m^{j}/l^{j}$ are both strictly decreasing in $j$ and $ \lim_{j \rightarrow 0} h^{j}/m^{j} = \infty$ and $ \lim_{j \rightarrow 1} m^{j}/l^{j} = 0$ 
\end{block}
\begin{itemize}
\item Lower indexed workers have a comparative advantage in high skill and higher indexed workers have a comparative advantage in low skill. \medskip
\item At the extremes these comparative advantages are strong enough that there will always be some workers choosing to supply high and low skills.
\end{itemize}
\end{frame}

\begin{frame}{Worker thresholds}
\begin{block}{Lemma 2}
For any ratios of wages $w_{H}/w_{M}$ and $w_{M}/w_{L}$, there exist $J^{*}(w_{H}/w_{M})$ and $J^{**}(w_{M}/w_{L})$ such that:
\begin{itemize}
\item $t_{h}^{j} = 1$ for all $j < J^{*}(w_{H}/w_{M})$
\item $t_{m}^{j} = 1$ for all $j \in ( J^{*}(w_{H}/w_{M}), J^{**}(w_{M}/w_{L}))$
\item $t_{l}^{j} = 1$ for all $j >  J^{**}(w_{M}/w_{L})$
\end{itemize}
$J^{*}(w_{H}/w_{M})$ and $J^{**}(w_{M}/w_{L})$ are both strictly increasing in their arguments. 
\end{block}
$J^{*}(w_{H}/w_{M})$ and $J^{**}(w_{M}/w_{L})$ are defined such that:
\[
\frac{m^{J^{*}(w_{H}/w_{M})}}{h^{J^{*}(w_{H}/w_{M})}}=\frac{w_{H}}{w_{M}} \text{ and } \frac{l^{J^{**}(w_{M}/w_{L})}}{m^{J^{**}(w_{M}/w_{L})}}=\frac{w_{M}}{w_{L}}
\]
\end{frame}

\begin{frame}{Relative supply of skill}
The supply to skills is given by: \medskip
\[
L = \int_{J^{**}(w_{M}/w_{L})}^{1} l^{j}dj ; M = \int_{J^{*}(w_{H}/w_{M})}^{J^{**}(w_{M}/w_{L})} m^{j}dj ; H = \int_{0}^{J^{*}(w_{H}/w_{M})} h^{j}dj
\]
\begin{itemize}
\item An increase in $w_{H}/w_{M}$, given $w_{M}/w_{L}$, increases $H/M$ given that $J^{*}(w_{H}/w_{M})$ is increasing in $w_{H}/w_{M}$. \medskip
\item An increase in $w_{M}/w_{L}$, given $w_{H}/w_{M}$, increases $M/L$ given that $J^{**}(w_{M}/w_{L})$ is increasing in $w_{M}/w_{L}$.
\end{itemize}
We have two upward sloping relative supply curves. 
\end{frame}

\begin{frame}{Equilibrium with endogenous supply of skill}
\begin{block}{Proposition 5}
In the model with endogenous supplies, there exists a unique equilibrium summarized by $(I_{L}, I_{H}$, $P_{L}, P_{M}, P_{H}$, $w_{L}, w_{M}, w_{H}$, $J^{*}(w_{H}/w_{M}), J^{**}(w_{M}/w_{L})$, $L,M,H)$.
\end{block}
Proof: To prove uniqueness of the equilibrium requires a little more work and is relegated to the Theoretical Appendix.
\end{frame}

\begin{frame}{Comparative statics and interpretation}
\begin{itemize}
\item Following the automation of routine tasks, the reallocation of skills across tasks resulted in wage and employment polarization. \medskip
\item In response to these changes, some workers will also change their skill supply away from medium skills. \medskip
\item If the more elastic margin is between medium and low skills, a significant fraction of workers previously supplying medium skills will now supply low skill. \medskip
\item This gives a richer explanation of observed wage and employment polarization.
\end{itemize}
\end{frame}

\subsection{4.7 Offshoring}

\begin{frame}{4. A Ricardian model of the labor market}
\begin{itemize}
\item[\textcolor{gray}{4.1}] \textcolor{gray}{Environment}
\item[\textcolor{gray}{4.2}] \textcolor{gray}{Equilibrium without machines}
\item[\textcolor{gray}{4.3}] \textcolor{gray}{Special cases}
\item[\textcolor{gray}{4.4}] \textcolor{gray}{Comparative statics}
\item[\textcolor{gray}{4.5}] \textcolor{gray}{Task replacing technologies}
\item[\textcolor{gray}{4.6}] \textcolor{gray}{Endogenous choice of skill supply}
\item[\textcolor{red}{4.7}] \textcolor{red}{Offshoring}
\item[\textcolor{gray}{4.8}] \textcolor{gray}{Directed technical change}
\end{itemize}
\end{frame}

\begin{frame}{Offshoring}
\begin{itemize}
\item Besides technical progress, international trade is believed to have a major impact on labor markets. \medskip
\item Instead of simply trading finished products, there has been a greater tendency to engage in trade in tasks through offshoring. \medskip
\item Our analysis of technological progress directly translates to offshoring of routine and codifiable tasks (Proposition 6). \medskip
\item The task-based model provides a rich and unified framework to analyse the impacts of labor supply, technological change, and trade on labor markets.
\end{itemize}
\end{frame}

\subsection{4.8 Directed technical change}

\begin{frame}{4. A Ricardian model of the labor market}
\begin{itemize}
\item[\textcolor{gray}{4.1}] \textcolor{gray}{Environment}
\item[\textcolor{gray}{4.2}] \textcolor{gray}{Equilibrium without machines}
\item[\textcolor{gray}{4.3}] \textcolor{gray}{Special cases}
\item[\textcolor{gray}{4.4}] \textcolor{gray}{Comparative statics}
\item[\textcolor{gray}{4.5}] \textcolor{gray}{Task replacing technologies}
\item[\textcolor{gray}{4.6}] \textcolor{gray}{Endogenous choice of skill supply}
\item[\textcolor{gray}{4.7}] \textcolor{gray}{Offshoring}
\item[\textcolor{red}{4.8}] \textcolor{red}{Directed technical change}
\end{itemize}
\end{frame}

\begin{frame}{Directed technical change}
\begin{itemize}
\item We have rewritten the canonical model by allowing for the endogenous allocation of skill groups across tasks. \medskip
\item We have further extended the model by also allowing for the endogenous allocation of workers across skill groups. \medskip
\item Another significant aspect of the economic environment absent from the canonical model is the endogeneity of technological progress to changes in supply. \medskip
\item This is known as endogenous or directed technical change.
\end{itemize}
\end{frame}

\begin{frame}{Directed technical change}
In our Ricardian model, there can be different types of directed technical change in response to changes in supply:
\begin{enumerate}
\item $A_{L}, A_{M}, A_{H}$ respond to changes in skill supplies. This idea is analyzed in Acemoglu \& Zilibotti [01] for the special case of our model discussed in section 4.3.
\item Comparative advantage schedules respond to changes in skill supplies:
\[
\alpha_{L}(i| \theta), \alpha_{M}(i| \theta), \alpha_{H}(i| \theta)
\]
with $\theta$ a variable that captures the endogenous choice of technology in the economy.
\end{enumerate}
\end{frame}

\begin{frame}{Factor biases in directed technical change}
\begin{itemize}
\item An increase in the supply of factor $f=L,M,H$ induces technical change that is \textbf{weakly biased} towards $f$ if:
\[
\frac{\partial w_{f}(E_{-f}, E_{f} | \theta )}{\partial \theta}\frac{d \theta}{d E_{f}} \geq 0
\]
\item An increase in the supply of factor $f=L,M,H$ induces technical change that is \textbf{strongly biased} towards $f$ if:
\[
\frac{d w_{f}(E_{-f}, E_{f} | \theta )}{d E_{f}} = \frac{\partial w_{f}(E_{-f}, E_{f} | \theta )}{\partial E_{f}} + \frac{\partial w_{f}(E_{-f}, E_{f} | \theta )}{\partial \theta}\frac{d \theta}{d E_{f}} > 0
\]
\end{itemize}
\end{frame}

\begin{frame}{Weakly biased directed technical change}
\begin{block}{Proposition 7}
Under regularity conditions (which ensure the existence of a locally isolated equilibrium), an increase in the supply of factor $f$ (for $f \in {L,M,H}$) will induce technical change \textbf{weakly biased} towards that factor.
\end{block}
\begin{itemize}
\item Proof: Acemoglu [07] \smallskip
\item This is a strong theoretical reason to expect that the increase in educational attainment has induced the development of technologies favoring skilled workers. \smallskip
\item It does not necessarily imply that the wage of skilled workers must have increased. 
\end{itemize}
\end{frame}

\begin{frame}{Strongly biased directed technical change}
\begin{block}{Proposition 8}
Under regularity conditions (which ensure the existence of a locally isolated equilibrium), an increase in the supply of factor $f$ (for $f \in {L,M,H}$) will induce technical change \textbf{strongly biased} towards that factor - thus increasing the wage of that factor - if and only if the aggregate production possibilities set of the economy is locally nonconvex in factor $f$ and technology $\theta$.
\end{block}
\begin{itemize}
\item Proof: Acemoglu [07] \smallskip
\item Local nonconvexity in $f$ and $ \theta$ is common in models of endogenous technological change. \smallskip
\item The endogenous technology (or ``long-run'') demand for skill is upward sloping.
\end{itemize}
\end{frame}

\section{5. Comparative Advantage and Wages: An Empirical Approach}

\begin{frame}{Relative returns to skill versus tasks}
\begin{itemize}
\item The model predicts that following automation, $w_{L}/w_{M} \uparrow $ and $w_{H}/w_{M} \uparrow $, even as workers reallocate across tasks. \medskip
\item However, changes in relative wages are not the same as changes in relative task prices (e.g. they can be different for tasks that are reallocated between skill groups). \medskip
\item Regarding automation, the price of automatable tasks is expected to declined relative to abstract and manual tasks. \medskip
\item But tasks and task prices are not directly observed in data.
\end{itemize}
\end{frame}

\begin{frame}{Occupations as proxies for tasks}
\begin{itemize}
\item Let $ \gamma_{sejk}^{A}, \gamma_{sejk}^{R}, \gamma_{sejk}^{S}$ be the employment shares of Abstract, Routine and Service occupations in  demographic group $sejk$ (gender, education, age, region) in 1959. \medskip
\item By construction, we have that:
\[
\gamma_{sejk}^{A} + \gamma_{sejk}^{R} + \gamma_{sejk}^{S} = 1 
\]
\item Assume that in 1959 these groups have self-selected into task specialities according to comparative advantage, taking as given overall skill supplies and task demands (reflecting also available technologies and trade opportunities).
\end{itemize}
\end{frame}

\begin{frame}{Wages, skill and tasks}
\begin{itemize}
\item Let $w_{sejkt} $ be the mean log wage of a demographic group in year $t$ and $\Delta w_{sejk\tau}$ the change in $w$ in decade $\tau$. \smallskip
\item Estimate the regression model:
\begin{align*}
\Delta w_{sejk\tau} & = \sum_{t} \beta_{t}^{A} \gamma_{sejk}^{A} 1[\tau=t] + \sum_{t} \beta_{t}^{S} \gamma_{sejk}^{S} 1[\tau=t] \\ 
& + \delta_{\tau} + \phi_{e} + \lambda_{j} + \pi_{k} + e_{sejk\tau} \tag{40} \label{eq40} 
\end{align*}
\item $\beta_{t}^{A}$ and $\beta_{t}^{S}$ are decade specific slopes on the initial occupation shares in wage changes by demographic group (relative to the omitted   routine task shares with effects absorbed by $ \delta_{\tau}$). \smallskip
\item We expect that $\beta_{t}^{A}$ and $\beta_{t}^{S}$ will rise, and that $\delta_{t}$ will decline.
\end{itemize}
\end{frame}

\newpage
\begin{center}
\begin{figure}
\includegraphics[height=8.5cm,width=\textwidth]{table10a.png}
\end{figure} 
\end{center}
\newpage

\newpage
\begin{center}
\begin{figure}
\includegraphics[height=8.5cm,width=\textwidth]{table10b.png}
\end{figure} 
\end{center}
\newpage

\section{6. Concluding Remarks}

\begin{frame}{Conclusions}
\begin{itemize}
\item Declining real wages for some workers and job and wage polarization are problematic for the canonical model. \smallskip  
\item A task-based framework, in which tasks are the basic unit of production and the allocation of skills to tasks is endogenous, provides a fruitful alternative theory. \smallskip 
\item The basic framework can be extended to include e.g. worker's choice of skill supply and directed technological change. \smallskip 
\item These models generate new ideas that can be tested empirically, but more needs to be done.
\end{itemize}
\end{frame}

\end{document}
